% Generated draft: requires mathematical and source review.
\documentclass[11pt]{article}
\usepackage[margin=25mm]{geometry}
\usepackage{fontspec}
\setmainfont{Latin Modern Roman}
\usepackage{amsmath,amssymb,mathtools,unicode-math}
\setmathfont{Latin Modern Math}
\usepackage{enumitem,needspace,xurl,hyperref}
\setlength{\parindent}{0pt}
\setlength{\parskip}{5pt}
\setlength{\emergencystretch}{3em}
\newcommand{\N}{\mathbb{N}}
\newcommand{\Z}{\mathbb{Z}}
\newcommand{\Q}{\mathbb{Q}}
\newcommand{\R}{\mathbb{R}}
\newcommand{\C}{\mathbb{C}}
\newcommand{\F}{\mathbb{F}}
\newcommand{\A}{\mathbb{A}}
\newcommand{\PP}{\mathbb{P}}
\newcommand{\E}{\mathbb{E}}
\newcommand{\eps}{\varepsilon}
\newcommand{\dd}{\,\mathrm{d}}
\begin{document}
\section*{UnsolvedMath ID: 700753}
\textbf{Source problem code:} catalogue entry
\subsection{Definitions and mathematical statement}
Let \(p\) be a prime and let \(F\) be a non-abelian free pro-\(p\) group, equipped with its pro-\(p\) topology. No finite-generation or finite-rank hypothesis is stated. For \(r\in F\), write
\[
(r)=\bigcap\{N\trianglelefteq F:N\text{ is closed and }r\in N\}
\]
for the closed normal subgroup generated by \(r\). The quotient \(F/(r)\) carries its quotient topology.

For a pro-\(p\) group \(G\), let \(\operatorname{cd}_p(G)\) denote its continuous \(p\)-cohomological dimension: the least integer \(d\geq0\) such that \(H^n_{\mathrm{cts}}(G,M)=0\) for every \(n>d\) and every discrete continuous \(G\)-module \(M\) whose underlying abelian group is \(p\)-primary torsion, or \(\infty\) if no such integer exists. The source writes \(\operatorname{cd}\) without specifying a coefficient convention; the statement below uses the standard pro-\(p\) interpretation \(\operatorname{cd}=\operatorname{cd}_p\). This interpretation requires confirmation in mathematical review.

The question attributed to O. V. Mel'nikov is whether, for every such \(p\) and \(F\), the set
\[
D_F=\{r\in F:\operatorname{cd}_p(F/(r))\leq2\}
\]
is dense in \(F\). Equivalently, is it true that for every \(x\in F\) and every open normal subgroup \(U\trianglelefteq F\), there exists \(r\in xU\) such that \(\operatorname{cd}_p(F/(r))\leq2\)?
\subsection{Short English statement}
For every non-abelian free pro-\(p\) group, must every nonempty open subset contain an element whose closed normal closure gives a quotient of cohomological dimension at most two? Here cohomological dimension is interpreted as continuous \(p\)-cohomological dimension, a convention not explicitly defined in the supplied source.
\subsection{Sources}
\begin{itemize}
\item Local UnsolvedMath registry, supplied complete record with permanent ID \(700753\), exact source code \texttt{catalogue entry}, and upstream revision \texttt{d371ee917ac9647fbb06cd6fc529deec6f3dda82}.
\item \emph{The Kourovka Notebook (Unsolved Problems in Group Theory)}, Problem 11.61, proposed by O. V. Mel'nikov; attributed in the supplied record to the 11th issue (1990), with the 21st issue edition (2026), \texttt{arXiv:1401.0300v48}, also identified. The collection description names E. I. Khukhro and V. D. Mazurov as editors.
\item The local record labels the problem \texttt{open} and records an update timestamp of \texttt{2026-10-16T14:00:00Z}. This is supplied registry metadata, not an independently verified assessment of the literature. The cited source was not independently consulted for this draft.
\end{itemize}
\end{document}
